\documentclass[11pt,a4paper]{article}
\usepackage[T1]{fontenc}
\usepackage[utf8]{inputenc}
\usepackage{amsmath,amssymb,xcolor}
\usepackage[margin=25mm]{geometry}
\usepackage[colorlinks=true,urlcolor=blue]{hyperref}
\definecolor{nexorablue}{HTML}{2459B5}
\setlength{\parindent}{0pt}
\setlength{\parskip}{7pt}
\setlength{\emergencystretch}{2em}
\renewcommand{\familydefault}{\sfdefault}
\begin{document}

{\small\bfseries\color{nexorablue} NEXORA CPGE}\par
{\small MP / NEX-SLC-06 / Énoncé}\par
{\LARGE\bfseries Choisir le plus petit gain conforme au cahier des charges}\par
\vspace{0.25em}\hrule\vspace{0.6em}

Un procédé $P(p)=1/(p+4)$ est commandé par un gain proportionnel $K\ge0$ avec retour unitaire négatif. On exige une constante de temps au plus égale à $0{,}10\,\mathrm s$ et un écart permanent au plus égal à $0{,}20$ pour une consigne unitaire. Déterminer le plus petit gain, puis vérifier les deux exigences.


\end{document}
